\chapter{Teoría cinética de los gases}
\label{cap:teoria_cinetica}

La función de distribución de rapidez de Maxwell-Boltzmann describe la distribución de rapideces de las moléculas en un gas:
$$ N_v = 4\pi N \left( \frac{m_0}{2\pi k_B T} \right)^{3/2} v^2 e^{-m_0 v^2 / 2k_B T} \quad (8.42) $$
$$ v_{rms} = \sqrt{\overline{v^2}} = \sqrt{\frac{3k_B T}{m_0}} = 1.73 \sqrt{\frac{k_B T}{m_0}} \quad (8.43) $$
$$ v_{prom} = \sqrt{\frac{8k_B T}{\pi m_0}} = 1.60 \sqrt{\frac{k_B T}{m_0}} \quad (8.44) $$
$$ v_{mp} = \sqrt{\frac{2k_B T}{m_0}} = 1.41 \sqrt{\frac{k_B T}{m_0}} \quad (8.45) $$

\section*{Piense, dialogue y comparta}
Consulte el Prefacio para obtener una explicación de los iconos utilizados en este conjunto de problemas.

\begin{enumerate}
    \subimport{piense_01/}{ejercicio.tex}
    \subimport{piense_02/}{ejercicio.tex}
    \subimport{piense_03/}{ejercicio.tex}
\end{enumerate}

\section*{Problemas}

\subsection*{SECCIÓN 8.1 Modelo molecular de un gas ideal}
\begin{enumerate}
    \subimport{prob_8_1_01/}{ejercicio.tex}
    \subimport{prob_8_1_02/}{ejercicio.tex}
    \subimport{prob_8_1_03/}{ejercicio.tex}
    \subimport{prob_8_1_04/}{ejercicio.tex}
    \subimport{prob_8_1_05/}{ejercicio.tex}
    \subimport{prob_8_1_06/}{ejercicio.tex}
    \subimport{prob_8_1_07/}{ejercicio.tex}
    \subimport{prob_8_1_08/}{ejercicio.tex}
\end{enumerate}

\subsection*{SECCIÓN 8.2 Calor específico molar de un gas ideal}
\begin{enumerate}
    \setcounter{enumi}{8}
    \subimport{prob_8_2_09/}{ejercicio.tex}
    \subimport{prob_8_2_10/}{ejercicio.tex}
    \subimport{prob_8_2_11/}{ejercicio.tex}
    \subimport{prob_8_2_12/}{ejercicio.tex}
    \subimport{prob_8_2_13/}{ejercicio.tex}
\end{enumerate}

\subsection*{SECCIÓN 8.3 Equiparición de la energía}
\begin{enumerate}
    \setcounter{enumi}{13}
    \subimport{prob_8_3_14/}{ejercicio.tex}
    \subimport{prob_8_3_15/}{ejercicio.tex}
    \subimport{prob_8_3_16/}{ejercicio.tex}
    \subimport{prob_8_3_17/}{ejercicio.tex}
\end{enumerate}

\subsection*{SECCIÓN 8.4 Procesos adiabáticos para un gas ideal}
\begin{enumerate}
    \setcounter{enumi}{17}
    \subimport{prob_8_4_18/}{ejercicio.tex}
    \subimport{prob_8_4_19/}{ejercicio.tex}
    \subimport{prob_8_4_20/}{ejercicio.tex}
    \subimport{prob_8_4_21/}{ejercicio.tex}
\end{enumerate}

\subsection*{SECCIÓN 8.5 Distribución de rapideces moleculares}
\begin{enumerate}
    \setcounter{enumi}{21}
    \subimport{prob_8_5_22/}{ejercicio.tex}
    \subimport{prob_8_5_23/}{ejercicio.tex}
    \subimport{prob_8_5_24/}{ejercicio.tex}
    \subimport{prob_8_5_25/}{ejercicio.tex}
    \subimport{prob_8_5_26/}{ejercicio.tex}
\end{enumerate}

\section*{PROBLEMAS ADICIONALES}
\begin{enumerate}
    \setcounter{enumi}{26}
    \subimport{prob_adicional_27/}{ejercicio.tex}
    \subimport{prob_adicional_28/}{ejercicio.tex}
    \subimport{prob_adicional_29/}{ejercicio.tex}
    \subimport{prob_adicional_30/}{ejercicio.tex}
    \subimport{prob_adicional_31/}{ejercicio.tex}
    \subimport{prob_adicional_32/}{ejercicio.tex}
    \subimport{prob_adicional_33/}{ejercicio.tex}
    \subimport{prob_adicional_34_a/}{ejercicio.tex}
    \subimport{prob_adicional_34_b/}{ejercicio.tex}
    \subimport{prob_adicional_35/}{ejercicio.tex}
    \subimport{prob_adicional_36/}{ejercicio.tex}
    \subimport{prob_adicional_37/}{ejercicio.tex}
    \subimport{prob_adicional_38/}{ejercicio.tex}
    \subimport{prob_adicional_39/}{ejercicio.tex}
    \subimport{prob_adicional_40/}{ejercicio.tex}
    \subimport{prob_adicional_41/}{ejercicio.tex}
    \subimport{prob_adicional_42/}{ejercicio.tex}
    \subimport{prob_adicional_43/}{ejercicio.tex}
    \subimport{prob_adicional_44/}{ejercicio.tex}
\end{enumerate}

\section*{PROBLEMAS DE DESAFÍO}
\begin{enumerate}
    \setcounter{enumi}{44}
    \subimport{prob_desafio_45/}{ejercicio.tex}
\end{enumerate}
